\section*{Summary}

The nuclear response of \cref{chap:11} hands a recoil over to dynamics on the Born--Oppenheimer surface through one exact step: the Wigner function of the lattice with the momentum of the struck nucleus raised by $\vect{q}$ (\cref{sec:19.1}). In silicon the dynamics passes from quantum to classical at \SI{0.41}{eV}, leaves the harmonic regime near \SI{0.54}{eV}, displaces atoms from 12.5 to \SI{36}{eV} depending on direction, and excites electrons nonadiabatically from about \SI{0.9}{keV}. Below the displacement threshold a recoil deposits all of its energy as heat, which covers every electron-following recoil from dark matter lighter than about \SI{150}{MeV}.

A binary collision with a valence electron can bridge the gap only above $E_{R,\mathrm{th}} = (M_\kappa/m_e)\Eg^2/16E_{\mathrm{F}}$, \SI{352}{eV} in silicon and \SI{399}{eV} in germanium, yet the measured silicon yield shows no threshold down to \SI{100}{eV} (\cref{sec:19.2}). The Migdal part of the measured yield is at most \SI{0.4}{\percent} of the silicon yield at \SI{100}{eV}. Lindhard quenching is the limit with no stored energy, continuum stopping to zero speed and proportional conversion to pairs (\cref{sec:19.3}); the silicon data lie 8 to \SI{38}{\percent} below the formula between 0.1 and \SI{3.9}{keV}. Bethe stopping holds within \SI{10}{\percent} only above $1.84\,v_{\mathrm{F}}$, while silicon recoils move at 0.013 to $0.13\,v_{\mathrm{F}}$.

One balance law per event counts every share of the deposit once (\cref{sec:19.5}). A close Frenkel pair in silicon keeps 6.8 to \SI{7.5}{eV}, and the displacement count with a screened cross-section is 0.70 times the Kinchin--Pease value, against the 0.8 of the standard formula. Unit heat quenching requires that, within the read window, all gap energy has returned as phonons and no energy is stored. Under conditions on phonon collection, other missing energy and the pair energy of the reconstruction, the measured deficits of nuclear recoils bound the stored share at \SI{95}{\percent} confidence by 0.055 in silicon and by 0.064 in germanium. The silicon datum $Y_{\mathrm{ph}} = 0.952$ is reproduced either by a stored fraction of 0.042 or, with no stored energy, by a collection ratio of 1.128; at \SI{40}{V} the two readings differ by 0.028.

The stored energy is released after the partition time by tunneling (\cref{sec:19.6}). The release-path sum rule \cref{eq:19.6-sumrule} makes the energy released by all hops equal to the stored energy minus what is still stored, for any rates. A single release gives a sensor that absorbs a quarter of the phonons less than \SI{2}{eV}, and two-level systems give quanta of half a meV, so neither produces the excess events of tens of eV. Energy conservation bounds the low-energy excess by the stored energy of the events that seeded it and of the grown-in defects, plus the decrease of relaxable elastic energy, and the factor of a hundred between glued and suspended crystals places it in the elastic energy. The uncertainty budget of the chapter is \cref{tab:19-nuisances}. The charge, phonons and quasiparticles that the balance law leaves at the end of the read window are not yet a signal; the next chapter follows them to the sensors.
